\section{The Complete CAD Battery}
\subsection{A Matrix of Different Evidence Kinds}
The historical battery contains 23 retained subjects and 18 evaluator columns. Its 414 cells include software capability, geometry binding, numerical findings, shared reports and record completeness. The manifest fixes membership. A later automatic assessment is a separate corpus and does not silently create a twenty-fourth subject. All cells in this historical corpus report \code{fresh_this_lane: false}.

The stored statuses reconcile to \BatteryPass\ passes, \BatteryFail\ failures and \BatteryBlocked\ blocked records. These are labels from heterogeneous producers. They should be preserved as historical data while the explanation states what each label meant. In particular, the four shared build columns account for 92 stored passes, and the local-solid column's favorable status concerns completeness. Neither fact supports an aggregate claim of 221 independently successful physical engineering tests.

% Generated stored-label counts, not physical-test outcomes.
\begin{table}[t]
\caption{Historical battery outcomes; P/F/B are stored labels.}
\centering\small
\begin{tabular}{lrrr}\toprule
Evaluator & P & F & B\\\midrule
build.analytic & 23 & 0 & 0\\
build.gci\_contracts & 23 & 0 & 0\\
build.claims & 23 & 0 & 0\\
build.dolfinx & 23 & 0 & 0\\
geom.unit & 0 & 0 & 23\\
geom.one\_solid & 0 & 0 & 23\\
geom.mass & 19 & 0 & 4\\
geom.pin & 0 & 0 & 23\\
plate.loads & 0 & 19 & 4\\
plate.seq & 10 & 9 & 4\\
plate.l7 & 0 & 19 & 4\\
plate.gci & 4 & 15 & 4\\
side.a1 & 19 & 0 & 4\\
side.a2 & 19 & 0 & 4\\
side.a3 & 19 & 0 & 4\\
side.a4 & 19 & 0 & 4\\
solid.tip & 20 & 0 & 3\\
gate.rs\_v2 & 0 & 19 & 4\\
\midrule
Total & 221 & 81 & 112\\\bottomrule
\end{tabular}
\label{tab:battery}
\end{table}


\test{build.analytic: Analytic Software Controls}
This column references shared analytic/reference software controls. Such controls can expose errors in units, invariance, assembly or solver behavior on reference problems. A per-subject copy of a shared control status does not establish another execution. The historical 23 passes should therefore be understood as repeated references to capability evidence. They do not establish candidate-specific geometry validity or physical behavior. The referenced control source and an execution receipt answer different questions and should not be conflated.

\test{build.gci\_contracts: Convergence Arithmetic}
This column concerns software rules for interpreting mesh-study data. Positive and negative controls can demonstrate that a formula or refusal condition is implemented as intended. That is valuable code-level evidence. It does not mean that a particular CAD has a suitable three-level study, matching observables, or an admitted asymptotic error estimate. Candidate-specific convergence belongs to the separate mesh-study evidence. Repeated software-control passes cannot fill missing numerical observations in that study.

\test{build.claims: Reporting and Admission}
This column references controls over how evidence is accepted, refused and rendered into claims. Its object is the reporting system. A useful negative control attempts to support an assertion with mismatched, incomplete or unauthorized evidence and verifies that the system refuses the claim. A favorable software outcome establishes behavior of that control under its fixture. It does not establish the mechanics of any CAD or authorize an unscoped public safety statement. Current source availability also does not prove that a historical control was freshly executed.

\test{build.dolfinx: Separate-Solver Capability}
The retained column references shared solver-comparison evidence. It indicates the scope of that shared capability record, not a fresh dual-solver computation for every subject. Candidate-specific comparison needs the two actual execution identities, matched problem specification and observable correspondence. If those are absent, the report must not infer them from the shared pass. The distinction prevents an installed solver or a historical demonstration from becoming an unsupported claim that every displayed candidate was independently checked.

\test{geom.unit: Export Units}
The unit column asks whether the exported representation's unit interpretation is established by the relevant evidence. All 23 historical cells are blocked. This does not prove that every export has incorrect units; it means the required binding was not established in those records. A plausible overall size is insufficient because an implicit unit error can preserve a convincing shape. A current importer or a later unit check cannot retroactively supply a missing historical receipt without a separately identified new observation.

\test{geom.one\_solid: Topology Binding}
This column records whether the relevant single-solid topology was established for the subject. Its 23 blocked cells preserve a missing-evidence condition rather than 23 unfavorable topology measurements. The distinction matters because topology can be assessed on different representations. A mesh closure check and a B-rep solid check do not establish the same fact. Reports should identify the representation and analyzer, preserve original status, and keep any later evidence separate from the historical cell.

\test{geom.mass: Mass Observation}
The mass column has 19 retained genome observations and four blocked reference entries. For those genome entries, the historical value is a plate/certification mass observation. It is not a demonstrated STEP-versus-raster consistency comparison. A mass number needs units, representation, density interpretation and computation origin. Agreement of rounded values from related pipelines does not establish identical volume or geometry. This column's stored success therefore means the observation is present under its historical rules, not that every relevant mass representation has been independently reconciled.

\test{geom.pin: Geometry Identity}
This column concerns binding the observation to the intended CAD. All historical entries are blocked. A content hash, when present and correctly associated, can identify bytes; a filename alone cannot. The model-to-artifact relationship also requires semantic information about export and conversion. A later source revision should not silently repair the historical record. Instead, the report should show the old missing binding and any separately retained new evidence, allowing readers to distinguish original observations from subsequent corrections.

\test{plate.loads: Plate Response Finding}
The load column reports the stored plate-response comparison. Nineteen records are unfavorable and four are blocked. The unfavorable findings are results of the model and comparison policy, not evidence that the computation itself failed to run. Conversely, blocked entries should not be described as measured failures. The plate formulation, field semantics, material envelope and source revision determine the interpretation. A report that preserves these distinctions communicates more engineering information than a single red/green matrix legend.

\test{plate.seq: Regional Response Comparison}
The sequence column records the retained regional comparison, with ten stored favorable and nine unfavorable findings among genome entries. It inherits the elastic-proxy limitations described for G16-14. The favorable label does not establish a physical fracture sequence. A reviewer needs the regions, governing case and ratio convention to compare observations meaningfully. This column and the overlapping grouped predicate do not supply independent corroboration when both use the same source response.

\test{plate.l7: Fatigue Screen}
The fatigue column records 19 unfavorable model assessments and four blocked entries. Fatigue interpretation depends on more than static nominal material strength: process, surface, mean stress, spectrum, history and uncertainty can affect the relevance of a model. The historical record is useful as a disclosed screen, but it is not a lifetime test. No service-life claim follows from the existence of a fatigue factor, and an unfavorable factor is not cured by a favorable aggregate weighted score.

\test{plate.gci: Candidate Mesh Study}
This column reports candidate-specific convergence evidence: four favorable, fifteen unfavorable and four blocked records. It must be interpreted separately from the software arithmetic controls. A valid study needs the same observable across a defensible mesh sequence. A plateau can make an apparent small difference misleading, and a divergent sequence cannot be rescued by reporting only its finest result. The stored column preserves its producer's finding; the paper does not relabel those historical outcomes as newly admitted convergence estimates.

\test{side.a1: Shared Side-Report Reference}
This row references a shared retained side report. Its historical status does not establish a distinct executed A1 experiment or a uniquely identified observation field. Nineteen favorable references and four blocked entries should therefore be read as report availability under the historical producer's convention. The useful next step for a reader is to inspect the referenced report and its actual fields. Expanding a single shared report into multiple row names must not inflate the count of independent engineering evidence.

\test{side.a2: Shared Side-Report Reference}
The A2 row similarly references the common side report. Its name should not be used to invent a separate physical mechanism or execution that is not established by the record. A favorable status can support the narrow statement that the shared report was referenced under that column's rules. Any stronger interpretation needs a distinct field definition, source selector and execution identity. This paper retains the row because it exists in the corpus while making its shared nature explicit.

\test{side.a3: Shared Side-Report Reference}
The A3 row has the same evidential limitation: no distinct executed test or separate observation is established solely by its historical cell. Treating A1, A2 and A3 as three independent successful side experiments would manufacture corroboration from presentation structure. Their relationship is therefore part of the atlas explanation. The correct evidence unit is the actual report or observation with its provenance, rather than the number of labels used to display it.

\test{side.a4: Support and Stability Report}
The A4 column records the alternate-support/stability report. Unlike the three legacy references, its technical scope is associated with the support representation and linear stability model. Nevertheless, report completeness, a favorable numerical finding and physical support calibration remain distinct. The record should preserve solver status and support assumptions, and should not imply that idealized compliant supports reproduce a measured interface. The column's 19 stored passes do not establish nineteen physical contact experiments.

\test{solid.tip: Solid-Result Completeness}
This is the clearest example of why a generic pass label is insufficient. Twenty stored passes mean that local-solid result records were complete under the column's rules. They do not mean twenty favorable strength comparisons. A complete record can contain a low modeled response, missing independent verification, or an unfavorable grouped requirement. The report should expose completeness and the numerical finding side by side, allowing a reader to value the available computation without mistaking its existence for success.

\test{gate.rs\_v2: Combined Grouped Finding}
The final column records the combined historical RS-v2 finding. It contains nineteen failures and four blocked reference entries, with zero recorded passes. The denominator must retain that distinction: it is inaccurate to say that all 23 completed and failed the same assessment. A combined unfavorable result should link its constituent groups so readers can identify reused observations, missing evidence and actual deficits. It should not erase useful numerical results merely because the aggregate comparison is unfavorable.

\subsection{A Better Reading of the Matrix}
The matrix is valuable because it reveals coverage and missingness across a fixed corpus. Its scientific value increases when each row/column pair is traceable to the original kind, solver, evidence path, source commit and freshness flag. It decreases when heterogeneous labels are reduced to a large percentage of successful tests. The atlas therefore explains each column and retains exact IDs, while the paper uses the matrix as evidence about retained assessment coverage rather than physical reliability.
